\section*{Chapter \ref{ch:iv2:pricing-and-numerics} --- Pricing and Numerics}
\markboth{\protect\omTocMark{}{Answers: Pricing and Numerics}}{\protect\omTocMark{}{Answers: Pricing and Numerics}}

\begin{solution}{iq:iv2:pricing-and-numerics:1}
The price is $C = S\Phi(d_1) - Ke^{-rT}\Phi(d_2)$, with $d_{1,2} = [\ln(S/K) + (r \pm \sigma^2/2)T]/(\sigma\sqrt T)$. The rule of thumb gives \$8, against an exact 7.97. At 30\% and three months, it gives \$6, against 5.98.

The rule of thumb for an at-the-money call is $0.4\,S\sigma\sqrt T$, where 0.4 stands for $1/\sqrt{2\pi}$.

\iqwrong{Forgetting the root of time, and scaling price with $T$. A move grows with $\sqrt T$, and so does the option's value.}
\end{solution}

\begin{solution}{iq:iv2:pricing-and-numerics:2}
They cost the same. The expected return does not appear in the formula, because the hedge removes all exposure to it. If a trader objects, say that the option can be replicated, and the cost of replicating depends only on how much the stock moves, and not on where it is expected to go.

Delta-hedging leaves a portfolio with no first-order stock risk. The cost of that hedge depends on volatility. So the same price holds under any drift. A simulation of the hedged book, with drifts of 20\% and 2\%, gives the same mean result of about zero.

The objection does have one correct part: a high expected return does make the unhedged call a better bet. It just does not change the price at which it can be replicated.

\iqwrong{``The stock with a higher return must have dearer calls.'' Holding the price of the stock fixed, the higher return is already in that price.}
\end{solution}

\begin{solution}{iq:iv2:pricing-and-numerics:3}
The call is worth about \$6.67. The risk-neutral chance of going up is 1/3. If the true chance is 80\%, the price does not change.

The call pays 20 or 0. Hold $\Delta$ shares and a loan $L$: $120\Delta - L = 20$ and $90\Delta - L = 0$. So $\Delta = 2/3$ and $L = 60$. The cost is $100 \times 2/3 - 60 = 6.67$.

For the probability, $120q + 90(1 - q) = 100$, which gives $q = 1/3$. And $20 \times 1/3$ is the same 6.67.

\iqwrong{Pricing with the 80\%. That gives \$16, and at that price you sell the call and replicate it for \$6.67, whichever way the stock goes.}
\end{solution}

\begin{solution}{iq:iv2:pricing-and-numerics:4}
The standard error is 15 cents, and you need 2.25 million paths for 1 cent.

It is $\sigma/\sqrt N = 15/100$. For 0.01, $N = (15/0.01)^2 = 2\,250\,000$. A tenth of the error takes a hundred times the paths, so people turn to variance reduction.

\iqwrong{Thinking that doubling the paths halves the error. It takes four times the paths.}
\end{solution}

\begin{solution}{iq:iv2:pricing-and-numerics:5}
It is worth \$2. This is true in any model, because on every path exactly one of the two is alive, and together they pay what the plain call pays. As the barrier falls toward zero, the \$6 tends to \$8.

That is the in-out parity\index{in-out parity}. It holds path by path, and so it does not rest on a model of the stock. It does need the two to watch the barrier in the same way.

\iqwrong{Using this when the two contracts differ on rebate or monitoring. A rebate paid at the touch has to be added separately.}
\end{solution}

\begin{solution}{iq:iv2:pricing-and-numerics:6}
It is worth about \$0.49. The two together pay \$1 for sure, so together they are worth $e^{-0.05} \approx 0.951$. And $0.951 - 0.46 = 0.491$. You used no volatility.

You need the stock not to finish exactly at 100 with positive probability, or at least the two contracts must cover that case between them.

\iqwrong{Saying the two must add to 1. That forgets the discounting, since \$1 is paid in a year.}
\end{solution}

\begin{solution}{iq:iv2:pricing-and-numerics:7}
It is cheaper, because an average moves less than a final price. Use a volatility of $\sigma/\sqrt 3 \approx 17.3\%$, and the rule of thumb gives about \$6.93, compared with about \$12 for the ordinary call.

The variance of the average of a Brownian motion over a period $T$ is $\sigma^2T/3$. The start of the year counts toward the average at once, and it adds to every later point, and the end of the year counts for little. So the option's volatility falls by a factor of $\sqrt 3$.

\iqwrong{Using half the volatility on the grounds that the average sits midway. It is the variance that matters, and the answer is 1/3 of the variance.}
\end{solution}

\begin{solution}{iq:iv2:pricing-and-numerics:8}
It has 231 nodes if it recombines, and 2\,097\,151 if it does not. So a recombining tree is workable and the other is not. Path-dependent options, such as a lookback or an average, lose the saving.

With $n$ steps, a recombining tree has $(n + 1)(n + 2)/2$ nodes, and a bushy tree has $2^{n+1} - 1$. For a path-independent option, it does not matter how you got to a price, so you only need one value per price.

\iqwrong{Thinking any barrier option loses the saving. A barrier is fine on a recombining tree, since a node can be knocked out. It is only slow to converge.}
\end{solution}

\begin{solution}{iq:iv2:pricing-and-numerics:9}
For $E[Z]$, each pair gives exactly zero, so the estimate has no error. For $E[Z^2]$, each pair gives the same number twice, so half your draws are wasted. The trick helps when the payoff is close to monotonic in $Z$.

The pair's values are negatively correlated if the payoff is monotonic, so the errors cancel. For an even payoff, they are perfectly positively correlated, and it costs you twice as much for no more information. For an at-the-money call, the variance falls by about 37\% at the same cost.

\iqwrong{Using it everywhere as a free gain. A straddle is nearly even in $Z$, so antithetics hardly help or even hurt it.}
\end{solution}

\begin{solution}{iq:iv2:pricing-and-numerics:10}
Take $X = Z_1$ and $Y = 0.6Z_1 + 0.8Z_2$. With those draws, $Y = 0.6 - 0.8 = -0.2$. The coefficients are the second row of a Cholesky factor\index{Cholesky factor} of the correlation matrix.

The variance of $Y$ is $0.36 + 0.64 = 1$, and its covariance with $X$ is 0.6. Generally, if $LL^\top = \Sigma$, then $LZ$ has covariance $\Sigma$.

\iqwrong{Using $Y = 0.6Z_1 + 0.4Z_2$. That has variance 0.52, and the wrong correlation. The two squares have to add to 1.}
\end{solution}

\begin{solution}{iq:iv2:pricing-and-numerics:11}
The forward is about 105.13. Your contract is worth about \$4.88 today. You need no volatility, because you can replicate a forward with a position that never changes.

To get the stock at 100 in a year, you buy it now and borrow. The forward price is $100e^{0.05}$. A contract with $K = 100$ pays $S_T - 100$, which is worth $S - Ke^{-rT} = 100 - 95.12$.

\iqwrong{Thinking that the forward is the market's forecast. It comes from replication, and it has no view of the stock.}
\end{solution}

\begin{solution}{iq:iv2:pricing-and-numerics:12}
The work goes up by about 8. Twice as many price points, and four times as many time steps, since stability needs $\Delta t$ to shrink with the square of the price step.

So going from 100 to 200 points costs 8 times as much. An implicit scheme has no such limit, and it solves a tridiagonal system in time linear in the number of points, so it is the usual choice.

\iqwrong{Saying a factor of 2. The stability limit is the hidden cost of an explicit scheme.}
\end{solution}

\begin{solution}{iq:iv2:pricing-and-numerics:13}
Use $X = -\ln(1 - U)/2$. From $U = 0.5$, that is about 0.3466. The method is inverse-transform sampling\index{inverse-transform sampling}: $F^{-1}(U)$ has the distribution $F$. It is not the usual route for a normal because $\Phi^{-1}$ has no closed form, although good numerical versions do exist.

For a normal, Box-Muller or the ziggurat is common in a plain simulation. Quasi-random sequences need the inverse, because each uniform must map to a single normal.

\iqwrong{Thinking that $U = 0.5$ gives the mean, which is 0.5 here. It gives the median.}
\end{solution}

\begin{solution}{iq:iv2:pricing-and-numerics:14}
The chance is about 63\%.

Under the risk-neutral measure at a zero rate, $\ln S_T = \ln S - \sigma^2T/2 + \sigma\sqrt T\,Z$. The drift has to be $-\sigma^2/2$ for the stock to be a martingale. So $S_T > K$ means $Z > [\ln(K/S) + \sigma^2T/2]/(\sigma\sqrt T)$, which is $-d_2$. Here $d_2 = [\ln(1.1) - 0.0225]/(0.3\sqrt{0.5}) \approx 0.343$. So $P = \Phi(0.343) \approx 0.634$.

\iqwrong{Forgetting the $-\sigma^2/2$, or taking the wrong sign. That gives $d_1$, which is about 0.56, and a chance of 71\%.}
\end{solution}

\begin{solution}{iq:iv2:pricing-and-numerics:15}
The call is $E[S_T\mathbf 1\{S_T > K\}] - KP(S_T > K) = S\Phi(d_1) - K\Phi(d_2)$. The trick is to complete the square, or to change to the stock measure. For the example: $78.17 - 63.43 = 14.75$.

The second term comes straight from the previous probability. For the first, write $S_T = Se^{\sigma\sqrt T Z - \sigma^2T/2}$, so
\[
E[S_T\mathbf 1\{Z > -d_2\}] = S\int_{-d_2}^\infty \frac{e^{-(z - \sigma\sqrt T)^2/2}}{\sqrt{2\pi}}\,dz = S\Phi(d_2 + \sigma\sqrt T) = S\Phi(d_1).
\]
The shift by $\sigma\sqrt T$ is a change of numeraire\index{change of numeraire}. In the measure where the stock is the unit of account, the chance of finishing above $K$ is $\Phi(d_1)$.

A simulation of the two terms at a million paths gives the same answer.

\iqwrong{Writing $E[S_T \mid S_T > K]$ as $S$. The conditional mean is much higher than the unconditional one, and that is the entire point.}
\end{solution}

\begin{solution}{iq:iv2:pricing-and-numerics:16}
The equation is $\partial_tV + \tfrac12\sigma^2S^2\partial_{SS}V + rS\partial_SV - rV = 0$. The expected return has been cancelled by the hedge. The theta must be $-\tfrac12\sigma^2S^2\Gamma \approx -3.97$ a year.

Hold $V - \Delta S$. By Ito, its change is $(\partial_tV + \tfrac12\sigma^2S^2\Gamma)\,dt + (\partial_SV - \Delta)\,dS$. Choosing $\Delta = \partial_SV$ removes the $dS$ term, which carries the drift. What remains has no risk, so it must earn the rate: $r(V - \Delta S)\,dt$. That gives the equation.

For the number, $\tfrac12 \times 0.04 \times 10\,000 \times 0.0198 = 3.96$, and it is 3.97 with more digits. That is about 1.6 cents a trading day. Gamma and theta are one thing seen in two ways.

\iqwrong{Treating theta as a free-standing Greek. In this model, at a zero rate, it is gamma with a different sign.}
\end{solution}

\begin{solution}{iq:iv2:pricing-and-numerics:17}
The function is $g(t) = e^{\sigma^2(T - t)}$, so the value is $10\,000e^{0.04} \approx 10\,408$.

Substituting $V = S^2g$ gives $S^2g' + \tfrac12\sigma^2S^2 \cdot 2g = 0$, so $g' = -\sigma^2g$. With $g(T) = 1$, you get the result.

The expectation agrees. $E[S_T^2] = S^2E[e^{2\sigma\sqrt T Z - \sigma^2T}] = S^2e^{2\sigma^2T - \sigma^2T} = S^2e^{\sigma^2T}$.

\iqwrong{Giving $S^2$. That treats $E[S_T^2]$ as $E[S_T]^2$. The excess is the variance of $S_T$, which is 408.}
\end{solution}

\begin{solution}{iq:iv2:pricing-and-numerics:18}
A call pays the stock less the strike if it finishes above the strike, and that is the asset option less $K$ cash options. Here the asset one is worth \$53.98, the cash one \$0.4602 per \$1, and $53.98 - 46.02 = 7.97$.

The asset one is worth $S\Phi(d_1)$, not $S\Phi(d_2)$. When the stock finishes above the strike, it is above its average. So the conditional mean is larger than $S$, which is why it is worth more than $S$ times the chance.

\iqwrong{Thinking that delta is the chance. See the same split: $\Phi(d_1)$ is a chance only under the stock measure.}
\end{solution}

\begin{solution}{iq:iv2:pricing-and-numerics:19}
The call is worth about \$7.98, against 7.97 for the lognormal. The normal model's implied lognormal volatilities are about 22.4\% at 80 and 18.3\% at 120. So it has a downward skew, simply because a fixed dollar volatility is a higher percentage volatility at a lower price.

At the money, the call is $\sigma_N\sqrt{T/(2\pi)} = 20 \times 0.3989$. This explains the 0.4 rule, which is exact here. At $K = 80$, the normal price is $20\Phi(1) + 20\varphi(1) \approx 21.67$. Solving for the lognormal volatility gives the skew.

\iqwrong{Thinking that a normal and a lognormal model are the same at the same volatility. At the money they are nearly so. Away from it, one of them has a skew in the other's terms.}
\end{solution}

\begin{solution}{iq:iv2:pricing-and-numerics:20}
The chance is about 32\%, which is $2(1 - \Phi(1)) \approx 0.317$. With 252 daily checks, it falls to about 30\%.

For each path that touches 120 and ends below it, reflecting after the touch gives a path ending above it, with equal probability. So the chance of touching is twice the chance of ending above. The final distance is one standard deviation.

Monitoring only daily misses touches within a day. It is like moving the barrier up by $0.5826$ of a daily standard deviation, about \$0.73, which gives $2(1 - \Phi(1.037)) \approx 0.300$.

\iqwrong{Giving the 16\% chance of ending above 120. A path can touch 120 and fall back.}
\end{solution}

\begin{solution}{iq:iv2:pricing-and-numerics:21}
The perpetual contract is worth $100/110 \approx \$0.909$. The one-year rule of thumb gives 0.564, and the exact value is 0.603.

Since $S$ is a martingale that goes to zero or hits 110, and since optional stopping applies when it is bounded by 110, $100 = 110p + 0$.

For one year, the European digital at 110 is $\Phi(d_2)$, with $d_2 = (-\ln 1.1 - 0.02)/0.2 = -0.577$, so 0.282. Twice that is 0.564. The exact formula adds a second term. With the drift $-\sigma^2/2$ on the log, $P = \Phi(-0.577) + (S/B)\Phi(-0.377) \approx 0.282 + 0.321 = 0.603$. So the rule of thumb is off because the log price has a drift. Reflection would make it exact for a walk with none.

\iqwrong{Using the reflection result for a lognormal stock without changing it. Reflection needs a symmetric walk, and the log price has a drift.}
\end{solution}

\begin{solution}{iq:iv2:pricing-and-numerics:22}
It is worth $-\sigma^2T/2 = -0.02$. It is negative because the log is concave, so a fair bet on the stock is a bad bet on its log. Selling the contract and holding $1/S$ shares gives you half of the realised variance, over a whole path.

By Ito, $d\ln S = dS/S - \tfrac12\sigma^2\,dt$. So $\ln(S_T/S_0) = \int dS/S - \tfrac12\int\sigma^2\,dt$. The first term is the gain on holding $1/S$ shares, with zero expected value. So that gain, less the log contract's payoff, is $\tfrac12\int\sigma^2\,dt$. That is the idea behind the variance swap\index{variance swap}, which holds for any path of volatility.

\iqwrong{Giving zero, by saying that the log's mean is the log of the mean. That is the error Jensen's inequality corrects.}
\end{solution}

\begin{solution}{iq:iv2:pricing-and-numerics:23}
It is worth about \$5.64, the same as a six-month at-the-money call today. It does not depend on the future stock level, because the value then is proportional to that level, and the expected level is today's.

At the start date, the call is worth $S_tf$, where $f$ is the value of a six-month at-the-money call on a stock worth 1. Under the risk-neutral measure, $E[S_t] = S_0$. So the value today is $S_0f = 5.64$.

\iqwrong{Forgetting what makes this work. The value on the start date scales with the level, and the level's expected value is today's. A volatility that depends on the level breaks the scaling, and with it the result.}
\end{solution}

\begin{solution}{iq:iv2:pricing-and-numerics:24}
It is worth about \$13.60, made up of 7.97 and 5.64. A one-year straddle would cost 15.93, so the chooser\index{chooser option} is cheaper, as you must decide early.

At the choice date, you take the larger of the call and the put. By parity with a zero rate, $P = C - S + K$. So $\max(C, P) = C + \max(K - S_t, 0)$. The second term is a put expiring at the choice date, and you hold the call for the whole time.

\iqwrong{Thinking a chooser is the same as a straddle. If you decide at the outset, you have a single option, and if you decide at the end, you have a straddle. A chooser is between them.}
\end{solution}

\begin{solution}{iq:iv2:pricing-and-numerics:25}
A negative price has a 2.3\% chance at each step. The Euler scheme gives $E[S_T^2]/S_0^2 = 1.25^4 \approx 2.44$. The true value is $e \approx 2.72$. Simulate the log instead, as this is exact for a constant volatility.

Each step multiplies by $1 + 0.5Z$. That is below zero if $Z < -2$, and $\Phi(-2) = 0.023$. For the second moment, $E[(1 + 0.5Z)^2] = 1.25$ per step. In the log scheme, $S_{i+1} = S_ie^{\sigma\sqrt{\Delta t}Z - \sigma^2\Delta t/2}$ cannot go negative, and it has the right distribution for any step size.

\iqwrong{Using more steps first. That does cure it, but slowly, and the exact scheme costs nothing extra. Euler is for when no exact scheme exists.}
\end{solution}

\begin{solution}{iq:iv2:pricing-and-numerics:26}
Use $u = e^{\sigma\sqrt{\Delta t}} \approx 1.105$, $d = 1/u \approx 0.905$ and $q = (1 - d)/(u - d) \approx 0.475$. It oscillates because the strike sits on a node for an even number of steps and between nodes for an odd number. Averaging two adjacent values cancels it: $(7.9457 + 7.9852)/2 = 7.9655$, which is very nearly exact.

The error falls like $1/n$, but it changes sign with the parity of $n$. Ten and eleven steps average to 7.958. So you can get three-digit accuracy from a tree too small to give it directly. For a strike off the node, the error also depends on where the strike falls between nodes.

\iqwrong{Trusting one tree price. Here each is off with a sign set by the parity of its steps, and 100 and 101 steps bracket the answer.}
\end{solution}

\begin{solution}{iq:iv2:pricing-and-numerics:27}
The largest stable step is $\Delta t = 1/1\,600$ of a year, so 1\,600 steps for a year. A little beyond that, the top of the grid develops a sawtooth and blows up.

The update is $V_j^{\text{new}} = V_j + a_j(V_{j+1} - 2V_j + V_{j-1})$, with $a_j = \tfrac12\sigma^2S_j^2\Delta t/\Delta S^2$. For stability, $a_j \leq \tfrac12$. That is tightest at $S = 200$: $\Delta t \leq 1/(0.04 \times 200^2)$.

The limit is set by the largest prices on the grid, far from anywhere you care about. An implicit scheme has no such limit, and it costs the same per step, up to a constant.

\iqwrong{Choosing the time step from the price level of interest, $S = 100$. That gives four times the limit, and the far edge blows up.}
\end{solution}

\begin{solution}{iq:iv2:pricing-and-numerics:28}
The Cholesky fails with a negative pivot, since no such three variables exist. The determinant is $-0.62$, and the smallest eigenvalue is about $-0.27$. The lowest allowed correlation of B and C is 0.62. In practice, you take the nearest valid matrix.

If two variables each correlate 0.9 with A, they cannot be unrelated. The bound is $\rho_{BC} \geq 0.9^2 - 0.19 = 0.62$, since $\sqrt{(1 - 0.81)(1 - 0.81)} = 0.19$.

To repair it, take the eigenvalues, clip the negative ones to a small positive value, and rescale so that the diagonal is 1. A better method is Higham's nearest correlation matrix\index{nearest correlation matrix}.

\iqwrong{Thinking that any symmetric matrix of numbers between $-1$ and 1 is a correlation matrix. This is a frequent bug in a risk system after someone overrides a single pair.}
\end{solution}

\begin{solution}{iq:iv2:pricing-and-numerics:29}
The rule of thumb gives $12/40 = 30\%$. One Newton step gets to 30.194\%, and a second to 30.1938\%, already enough for six digits.

The step is $\sigma_{\text{new}} = \sigma - (C(\sigma) - 12)/\text{vega}$. The price at 30\% is 11.9235, and the vega is 39.4. So the step is $+0.00194$. It is so quick because an at-the-money price is nearly linear in volatility.

It fails where the vega is small: deep in or out of the money, or close to expiry. There it can overshoot to a negative volatility. So you keep a bracket and fall back on bisection, or Brent, and first check that the price lies within the no-arbitrage bounds.

\iqwrong{Skipping the bounds check. No volatility gives a price below intrinsic value, and the iteration then runs forever on a bad quote.}
\end{solution}

\begin{solution}{iq:iv2:pricing-and-numerics:30}
With different random numbers, the standard error is about 0.029. With the same ones, it is about 0.002, which is about 16 times less. That is 250 times less variance, for nothing.

With independent runs, the estimate is $(A - B)/2$, where each has an error of $13.15/\sqrt{100\,000} = 0.0416$. So the error is $\sqrt 2 \times 0.0416/2 = 0.029$. That is useless if delta is about 0.54.

With the same draws, the two payoffs nearly cancel path by path. What remains is a value of 0 or about 1 on each path. Its standard deviation is about 0.59, and $0.59/\sqrt{100\,000} \approx 0.0019$.

\iqwrong{Making the bump smaller to get precision. With independent runs, the error grows as the bump falls, in proportion to $1/h$.}
\end{solution}

\begin{solution}{iq:iv2:pricing-and-numerics:31}
The basket's volatility is about 26\%, and the single-lognormal price is about \$10.34. A simulation gives about 10.37. With a correlation of 0, it is about \$8.45, and with a correlation of 1, about \$11.92.

The volatility: $0.25 \times 0.09 \times 2 + 2 \times 0.25 \times 0.5 \times 0.09 = 0.0675$, whose root is 0.2598. Matching the true second moment, $E[B^2]/10\,000 = \tfrac12e^{0.09} + \tfrac12e^{0.045}$, gives 26.03\% and about 10.35. The basket is a sum of lognormals and not itself lognormal, so the approximation is close but not exact, and a simulation is the check.

\iqwrong{Averaging the two volatilities, 30\%. That is the correlation-of-1 case, and the diversification is what you are being paid for.}
\end{solution}

\begin{solution}{iq:iv2:pricing-and-numerics:32}
The variance falls by the fraction $\rho^2$, so about 79\%. So what remains is about 21\%, and the error is about 46\% of what it was, or the same as using about 4.8 times the paths. The optimal multiplier is $\text{Cov}(\text{payoff}, S_T)/\text{Var}(S_T)$, about 0.58. You estimate it by a regression on a pilot sample.

The new estimate is $\bar Y - \beta(\bar X - E[X])$. Its variance is $\text{Var}(Y)(1 - \rho^2)/N$. So an at-the-money call has a standard deviation of 13.15 per path before, and about 6.0 after.

A deep in-the-money call nearly equals the stock, so its correlation is close to 1, and the gain is enormous. A deep out-of-the-money call behaves like a rare event, so the correlation falls, and the gain shrinks.

\iqwrong{Fitting $\beta$ on the same paths and saying nothing about it. That leaves a small bias, and with few paths you should use a pilot sample.}
\end{solution}

\begin{solution}{iq:iv2:pricing-and-numerics:33}
The volatility is $\sigma/\sqrt 3 \approx 17.3\%$, the forward is $Se^{-\sigma^2T/12} \approx 99.25$, and the price is about \$6.51. There is a closed form since a sum of normals is normal, while a sum of lognormals has no simple law.

The log of the average is $\ln S + \frac1T\int_0^T(\sigma W_t - \tfrac12\sigma^2t)\,dt$. It is normal. Its mean is $\ln S - \sigma^2T/4$, and the variance of the integral is $\sigma^2T/3$. So the forward is $S\exp(-\sigma^2T/4 + \sigma^2T/6) = Se^{-\sigma^2T/12}$.

Apply Black's formula to a forward of 99.25, a volatility of 17.3\% and a strike of 100. On 252 daily closes, the discrete version has a slightly higher volatility and is about 6.53.

\iqwrong{Using $S$ as the forward. A geometric average is below an arithmetic one, and its mean falls with variance. Here it matters by about 0.4.}
\iqfollow{What are the ways to price the arithmetic option?}{Three are common: a matched lognormal, which is quick but gets the tails wrong at high volatility; a simulation with this geometric option as a control; and a partial differential equation with the average as an extra variable, which a change of variable can bring back to one dimension.}
\end{solution}

\begin{solution}{iq:iv2:pricing-and-numerics:34}
The moment-matched price is about \$6.93. A simulation with the geometric option as a control variate confirms this at about 6.92 for daily averaging, with few paths. The standard error falls by about 25 times.

The first moment is $E[A] = S$ at a zero rate. The second gives $E[A^2] = 10\,306.9$. So the matched volatility is $\sqrt{\ln(1.03069)} \approx 17.39\%$.

For the check, simulate both payoffs with the same paths. The known geometric price is the control: $\bar P_A - \beta(\bar P_G - G)$. The residual standard deviation is $\sqrt{1 - 0.9992^2} \approx 0.04$ of the original, so a factor of 25. So 100\,000 paths give a standard error of 0.0014, where plain paths would give 0.036.

\iqwrong{Using the geometric price as the answer. Its forward is lower, and so it is too low by about 0.4. It is a control, and not an estimate.}
\iqwhy{Finding a related problem with a known answer is the most useful skill in simulation, and a desk uses this one daily.}
\end{solution}

\begin{solution}{iq:iv2:pricing-and-numerics:35}
It is worth about 10.52. The volatility of the ratio is $\sqrt{0.09 + 0.04 - 2 \times 0.5 \times 0.3 \times 0.2} = 26.5\%$. The rate does not enter because neither asset is being financed: this is Margrabe's formula\index{Margrabe's formula}.

The trick is to measure everything in units of B. Then the contract pays $(S_A/S_B - 1)^+$ in shares of B, and the ratio is a martingale. It is a call with strike 1 at the ratio's volatility, and with a zero rate, since B itself is the unit. So $V = S_B[\,(S_A/S_B)\Phi(d_1) - \Phi(d_2)]$, and at equal prices this is $100(2\Phi(0.1323) - 1)$.

As B's volatility goes to zero, the formula becomes Black's with $F = S_A$, at A's volatility, and with $K = S_B$. It is a call on A struck at a level that grows with the rate. The price also rises as correlation falls: with no correlation, the ratio's volatility is 36\% and the option is worth about 14.3.

\iqwrong{Using the volatilities' difference, 10\%. That is the case of perfect correlation, and it is the lowest possible.}
\iqfollow{How do you hedge it?}{Hold $\Phi(d_1) \approx 0.553$ of A and short $\Phi(d_2) \approx 0.447$ of B. Those two positions finance each other exactly as the prices move, so no cash account is needed.}
\end{solution}

\begin{solution}{iq:iv2:pricing-and-numerics:36}
It is worth exactly \$10. The volatility does not come in. An ordinary call at 30\% is worth about \$18.14, so the \$8.14 gap is the value of the knock-in part.

As long as the stock has not hit 100, the option pays $S_T - K$. On a path that hits the barrier, $S - K$ is zero there. So the payoff is the stopped value of the martingale $S - K$, which is worth $S - K$ today: optional stopping again.

The down-and-in call comes alive at the strike, when it holds nothing but time value, and that is the \$8.14. A dealer can replicate the down-and-out with one share and a loan of $K$, both closed at the barrier. This needs no gaps at the barrier, and the rate must be zero, or else equal to the carry.

\iqwrong{Trying to compute with the full barrier formula. It works, but it hides the point that no volatility enters. An interviewer asks this to see if you notice.}
\iqfollow{What if the barrier is at 95, below the strike?}{The volatility comes back, since a path can touch 100 and finish above it, and that is no longer the end of the matter. A rise in volatility now raises the value of a call that may still be alive at 100.}
\end{solution}

\begin{solution}{iq:iv2:pricing-and-numerics:37}
The prices are about 2.97 at 20\%, 3.87 at 15\% and 2.08 at 25\%. The vega is negative here, unlike an ordinary call, because a rise in volatility makes the knockout likelier.

The image formula is $V(S) = U(S) - (S/B)U(B^2/S)$, where $U$ is a European contract with the payoff $(S_T - K)^+$ cut off at $B$. Thus $U(x) = C(x, K) - C(x, B) - (B - K)D(x, B)$, with $D$ the digital. This $V$ meets the equation, and it is zero at the barrier, and its payoff is right. The factor $S/B$ comes from the drift of the log, with a zero rate.

An ordinary call only gains from volatility. Here the ordinary part is 7.97, and the extra probability of knocking out while deep in the money takes away more than the extra volatility adds.

\iqwrong{Assuming the vega is positive, since it is an option. The hedger's risk here is also unusual: near the barrier the option's value can fall to zero at once, and the delta goes negative.}
\end{solution}

\begin{solution}{iq:iv2:pricing-and-numerics:38}
It should be higher, because fewer checks mean fewer chances to die. The Broadie-Glasserman-Kou\index{Broadie-Glasserman-Kou shift} shift moves the barrier away by a factor $e^{0.5826\sigma\sqrt{\Delta t}}$, to about 130.96. The continuous formula at that barrier gives about 3.16. A simulation of daily checks gives about 3.14.

So a barrier move of under a dollar, in the far tail, changes the price by about 6\%. That is why you must always ask how a barrier is monitored.

The constant 0.5826 is $-\zeta(1/2)/\sqrt{2\pi}$. On a daily check, $\sigma\sqrt{\Delta t} = 0.0126$, so the factor is $e^{0.00734}$.

\iqwrong{Pricing a daily barrier by a continuous formula. The error is sizeable, and it is largest where the option has most of its value near the barrier.}
\end{solution}

\begin{solution}{iq:iv2:pricing-and-numerics:39}
The bias is that a path can cross 130 and come back between two monthly points. The fix is a Brownian-bridge\index{Brownian bridge} correction: you weight each step by the chance of having stayed below. That gives 2.97 at the same monthly steps. For this month, the chance is about 15\%.

Given two points, the path between them is a Brownian bridge in the log. The chance of crossing $b$ is $\exp[-2(b - x_1)(b - x_2)/(\sigma^2\Delta t)]$, in logs. Here $b - x_1 = \ln(130/120) = 0.0800$ and $b - x_2 = \ln(130/125) = 0.0392$, and $\sigma^2\Delta t = 0.04/12$. So the exponent is $-1.884$, and the chance is 0.152.

Multiply each path's payoff by the product of the survival chances over its steps. That has lower variance than sampling a crossing, and it cures the bias almost completely.

\iqwrong{Solving it with daily steps, thinking you have removed the bias. You have reduced it, as a daily simulation's 3.14 against 2.97 shows, but you have not removed it.}
\end{solution}

\begin{solution}{iq:iv2:pricing-and-numerics:40}
The pathwise estimate is $\mathbf 1\{S_T > K\}S_T/S_0$, with a mean of 0.54 and a standard deviation of 0.59. The likelihood-ratio estimate is $(S_T - K)^+Z/(S_0\sigma\sqrt T)$, with a mean of 0.54 and a standard deviation of 1.34. For a digital you must use the likelihood ratio, because its payoff is not differentiable.

Pathwise: differentiate the payoff along the path. Since $S_T$ is proportional to $S_0$, $\partial_{S_0}(S_T - K)^+ = \mathbf 1\{S_T > K\}S_T/S_0$.

Likelihood ratio: differentiate the density of the log price. The score is $Z/(S_0\sigma\sqrt T)$. It leaves the payoff alone, so it works with anything. But it has about 5 times the variance here, and much more for short-dated options.

For a digital, the pathwise derivative is zero almost everywhere, and the estimate is wrong: it gives zero. The other way is to smooth the digital into a tight call spread.

\iqwrong{Using pathwise on a digital and getting zero without noticing. A test of a Greek should compare with a closed form or with bumping.}
\end{solution}

\begin{solution}{iq:iv2:pricing-and-numerics:41}
It is worth about \$15.96, by reflection, and this is exactly the at-the-money straddle's value in this model. With daily checks, expect about 15.2.

By reflection, $\max W_T$ has the same law as $|W_T|$. So the mean is $\sigma\sqrt{2T/\pi} = 20 \times 0.798$. The straddle pays $|W_T|$ in the same units, so it has the same mean.

Daily checks lower the maximum by about $0.5826\sigma\sqrt{\Delta t} = 0.73$, and a simulation gives about 15.2. The exact mean with daily checks comes from Spitzer's identity: $E[\max] = \sum_k E[S_k^+]/k$, which is 15.24. The two contracts are not the same trade. The lookback's pay is the highest level, and the straddle's is the final distance from the start. They just have the same price. With a drift, or on a lognormal stock, the equality fails, because reflection needs a symmetric walk.

\iqwrong{Thinking the lookback must be worth twice the straddle, since it is ``perfect hindsight''. It only gets the highest level, and the straddle gets either side.}
\end{solution}

\begin{solution}{iq:iv2:pricing-and-numerics:42}
With no skew, it is worth 0.460. With the skew, it is about 0.540, up by about 8 cents. Most equity indexes have a downward smile, so this is the sign you usually see.

The digital is $-\partial_KC(K, \sigma(K))$, which is $-\partial_KC - \text{vega}\,\partial_K\sigma$. The first term is $\Phi(d_2)$. The second is $39.7 \times 0.002 = 0.079$.

The correction is positive with a downward skew, since a higher strike means lower implied volatility and a lower call value. A call spread built from real quotes gets this automatically. So use the spread, and not a formula with one volatility.

\iqwrong{Pricing a digital by plugging the strike's implied volatility into $\Phi(d_2)$. The error here is 8 cents on a 46-cent option, and it is the classic error on a digital.}
\end{solution}

\begin{solution}{iq:iv2:pricing-and-numerics:43}
The chance is about $3.17 \times 10^{-5}$, and the plain relative standard error is about 18\%. With the shift, the weight is $e^{-4x + 8}$ on a draw $x$, and the relative error falls to about 0.21\%.

Plain: $1/\sqrt{Np}$, with $Np = 31.7$ hits. Shifted: draw $X = 4 + Z$. The density ratio is $\varphi(x)/\varphi(x - 4) = e^{-4x + 8}$. About half of the draws hit, and they carry small, similar weights. The variance falls by a factor of about 7\,000.

In many dimensions, a shift can leave weights that are dominated by a few paths, and an estimated standard error can look fine and yet be wrong, if the sample has missed the region that matters. So watch the effective sample size, and think hard about where the rare event really lives.

\iqwrong{Shifting too far. Past the threshold, the variance rises again, since the weights become very unequal among the draws.}
\end{solution}

\begin{solution}{iq:iv2:pricing-and-numerics:44}
In this setting, plain Monte Carlo has an error of about 0.097, and a scrambled Sobol\index{Sobol sequence} about 0.0006, or 150 times less. The Sobol points fill the range evenly, so their errors cancel.

Quasi-Monte Carlo gets near a rate of $1/N$ on a smooth low-dimensional problem, in place of $1/\sqrt N$. Its advantage shrinks with many dimensions and with a discontinuous payoff. It also needs care with how the path's dimensions are ordered, so the Brownian bridge puts most of the variance into the first few coordinates.

Scramble the sequence, since then you can get an error bar from repeated seeds.

\iqwrong{Using one Sobol run and quoting $\sigma/\sqrt N$ as its error. Plain Monte Carlo's error bar does not apply, and only repeated scrambles give a real one.}
\end{solution}

\begin{solution}{iq:iv2:pricing-and-numerics:45}
The extrapolation is the best here, with an error of 0.0001. The usual textbook choice is the European control, and here it is worse than the plain tree: its error is $+0.0120$, against $-0.0080$. The average has an error of $+0.0031$. You exercise today below a level between \$81 and \$82.

Plain: $6.0824 - 6.0903$. Average: $(6.0824 + 6.1046)/2 - 6.0903$. Control: $6.0824 - 5.5536 + 5.5735 - 6.0903$. Extrapolation: $2 \times 6.0864 - 6.0824 - 6.0903$.

The control assumes the tree's error for the American and the European are similar. Here, the European's error is more than twice the American's. So do not assume it, and test the correction on your own problem against a very large tree. The extrapolation needs steps of the same parity, as the error's constant differs for odd and even.

\iqwrong{Quoting the control variate as a sure gain. The numbers here show that it is not one.}
\end{solution}

\begin{solution}{iq:iv2:pricing-and-numerics:46}
The scheme is stable, but it does not damp the high-frequency error. The kink in the payoff has high frequencies, and they survive as a sawtooth. So the prices look fine, and the gamma, a second difference, is garbage. The standard fix is Rannacher\index{Rannacher time-stepping} time-stepping: take the first few steps fully implicit. With two such steps, the gamma is 0.141.

Crank-Nicolson's amplification factor for a high-frequency mode is $(1 - x)/(1 + x)$, where $x$ is large when $\Delta t$ is big compared with $\Delta S^2$. That tends to $-1$, so the mode stays the same in size and flips sign at each step. An implicit step has the factor $1/(1 + 2x)$, which kills the mode.

So on this grid, near the strike, the second differences alternate: $0.87$, $-1.36$, $0.87$. After two implicit steps, they all lie near 0.141.

\iqwrong{Saying unconditional stability means accuracy. It only means the error does not grow. Here it does not shrink either.}
\end{solution}

\begin{solution}{iq:iv2:pricing-and-numerics:47}
The standard deviations are about 13.15 at strike 100, and about 4.82 at strike 130. So you need about 1.73 million paths and about 233\,000 paths for a 1-cent standard error. But as a share of the price, the error per path is 1.65 at 100 and 4.78 at 130.

The second moment is $S^2e^{\sigma^2T}\Phi(d_1 + \sigma\sqrt T) - 2SK\Phi(d_1) + K^2\Phi(d_2)$. It follows from the same change of measure as the price, used once for $S_T^2$ and once for $S_T$. The prices are 7.97 and 1.01.

So for a tail option, the dollar error is smaller, but the relative error is worse by about 3 times here. As a strike moves farther out, it worsens quickly, which leads to importance sampling.

\iqwrong{Quoting a dollar standard error for a cheap option. One cent on a \$1 option is 1\%, and on an option worth 10 cents, it is 10\%.}
\end{solution}

\begin{solution}{iq:iv2:pricing-and-numerics:48}
Simulate paths forward, then go backward. At each date, regress the discounted future cash on functions of the stock, using in-the-money paths only. Exercise where the immediate payoff beats the fitted value. With 100\,000 paths, you get about 6.04, and on fresh paths about 6.03, each with a standard error of 0.023. With 400 paths, the fit overstates: over many runs, in-sample estimates average about 6.2 and fresh-path ones about 5.9. One run has a plain standard error near 0.37, or a spread of about 0.2 with antithetic pairs, so it cannot show this on its own.

In steps:
\begin{enumerate}
\item Simulate paths at the 12 dates. Set each path's cash flow to its payoff at expiry.
\item Step back a date, and discount the cash flow.
\item For paths in the money, regress that discounted cash on a cubic in $S$.
\item Where $K - S$ is above the fitted value, replace the cash flow by $K - S$, at this date.
\item Repeat, then average the discounted cash at time 0.
\end{enumerate}

You use in-the-money paths only for two reasons. Out of the money, exercising is never right, so those paths cannot inform the choice. And including them spoils the fit in the region that matters.

On fresh paths, any rule you use is a valid exercise rule, and none beats the optimal rule. So the estimate is low-biased, up to noise. On the fitting paths, the rule has seen each path's future, which adds an upward bias. This bias grows as paths fall.

\iqwrong{Treating the regression's fitted value as the cash flow. Its role is only to decide. The estimate uses realised cash flows.}
\end{solution}

\begin{solution}{iq:iv2:pricing-and-numerics:49}
At every point, $V \geq K - S$, and $\partial_tV + \mathcal LV \leq 0$, and at least one of these is an equality. This is a linear complementarity problem. At each step, you solve it with projected successive over-relaxation\index{projected SOR}. With a \$1 grid and 100 steps, you get about 6.085, against 6.0903. The boundary today is between \$81 and \$82.

Here $\mathcal LV = \tfrac12\sigma^2S^2V_{SS} + rSV_S - rV$. So either you hold, and the equation holds, or you exercise, and the value is the payoff.

Each step is $Ax = b$, with $x \geq g$, where $g$ is the payoff. Projected SOR sweeps through the points, and after each update, it takes the larger of the new value and $g$. So the result satisfies all three conditions together.

If you solve the system and then cut it off at the payoff, you are explicit in the exercise decision: the point next to the boundary used a neighbour that was not yet corrected. That has lower accuracy in time, and it can break down at larger time steps with Crank-Nicolson. For a put, the Brennan-Schwartz algorithm is a direct and correct alternative, because the boundary has only one crossing.

\iqwrong{Forgetting to take a few fully implicit steps at the start (Rannacher time-stepping). With Crank-Nicolson, the kink in the payoff makes the solution oscillate.}
\end{solution}

\begin{solution}{iq:iv2:pricing-and-numerics:50}
The call is $C = SP_1 - Ke^{-rT}P_2$, where
\[
P_2 = \tfrac12 + \frac1\pi\int_0^\infty\operatorname{Re}\frac{e^{-iu\ln K}\phi(u)}{iu}\,du, \quad
P_1 = \tfrac12 + \frac1\pi\int_0^\infty\operatorname{Re}\frac{e^{-iu\ln K}\phi(u - i)}{iu\,\phi(-i)}\,du.
\]
For the test, at 100 it gives 7.965567, with $P_1 = 0.539828$ and $P_2 = 0.460172$. At 120 it gives 2.147299. At an upper limit of 200 in $u$, the loss is under $10^{-7}$.

$P_2$ is the chance of finishing above $K$, and $P_1$ is that same chance in the stock measure, where the characteristic function is $\phi(u - i)/\phi(-i)$. This is the Gil-Pelaez\index{Gil-Pelaez inversion} inversion formula.

In Black-Scholes, $\phi(u) = \exp(iu(\ln S - \sigma^2T/2) - \sigma^2u^2T/2)$. The tail in $u$ decays like a Gaussian with the scale $1/(\sigma\sqrt T)$, so a short-dated option needs a larger limit.

You would never do this for Black-Scholes. You test the code on it, as the exact answer is known. Then it applies to a model with no closed-form density, such as Heston or a jump-diffusion, whose $\phi$ is explicit.

\iqwrong{Taking the integral from $u = 0$ without care. The integrand has a finite limit there, but evaluating it directly gives $0/0$.}
\end{solution}
